\chapter{Conclusões}
\label{cap:conclusão}

% TODO: Parágrafo de abertura — resumir o trabalho em duas a três frases, com foco
% no que foi construído do ponto de vista de engenharia de software: uma API de
% backend robusta, uma estratégia de testes estruturada e uma infraestrutura de
% produção reprodutível e segura.

%----------------------------------------------------------------------------------------
\section{Objetivos Concretizados}
\label{sec:objetivos_concretizados}

% Retomar os três objetivos da introdução e avaliar o grau de concretização.
% Sugestão: lista estruturada ou tabela com "Objetivo" e "Resultado".
%
% 1. DESENVOLVIMENTO DO SERVIÇO CENTRAL
%    Concretizado: módulo de administração implementado e em produção, com gestão de
%    instituições, contas de caregiver e senior, controlo de acesso por perfil (RBAC)
%    e isolamento de dados por instituição. Fluxos de autenticação e onboarding
%    validados com utilizadores reais. O backend garante que o microserviço de IA
%    nunca é exposto diretamente, centralizando a autenticação e a construção
%    do contexto de cada conversa.
%
% 2. ESTRATÉGIA DE TESTES AUTOMATIZADOS
%    Concretizado: pirâmide de testes implementada com quatro níveis — unitários,
%    de contrato, de integração e de configuração de deploy. Cobertura dos serviços
%    de backend desenvolvidos durante o estágio com mocks das dependências externas.
%    Limitação: cobertura dos fluxos que envolvem respostas em streaming do microserviço
%    de IA ficou abaixo do desejado pela complexidade de simular respostas assíncronas
%    em ambiente de teste.
%
% 3. PROCESSO DE ENTREGA EM PRODUÇÃO
%    Concretizado: infraestrutura de produção operacional em AWS EC2 com contentores
%    Docker, proxy Nginx com TLS, estratégia blue-green para atualizações sem downtime
%    e migrações de base de dados automatizadas no arranque do contentor.
%    O processo de deploy é declarativo, versionado e reprodutível por qualquer
%    elemento da equipa.

%----------------------------------------------------------------------------------------
\section{Limitações e Trabalho Futuro}
\label{sec:trabalho_futuro}

% LIMITAÇÕES:
%   - Ausência de pipeline de CI/CD automatizado: o processo de deploy continua
%     semi-manual. As etapas de build, teste e promoção entre ambientes blue-green
%     são executadas manualmente, o que aumenta o risco de erro humano.
%   - Monitorização em produção: não existe ainda observabilidade estruturada —
%     logs centralizados, métricas de latência por endpoint e alertas automáticos
%     em caso de falha de um contentor ou fornecedor externo.
%   - Testes de carga: a infraestrutura não foi validada sob carga simultânea de
%     múltiplas conversas de voz, o que representa um risco desconhecido à medida
%     que o número de instituições cresce.
%
% TRABALHO FUTURO:
%   - Implementar um pipeline de CI/CD (e.g. GitHub Actions) que execute
%     automaticamente os testes, construa as imagens Docker e promova o deploy
%     após validação — eliminando a intervenção manual no processo de entrega.
%   - Adicionar observabilidade: logs estruturados (e.g. com Pino) exportados para
%     um agregador, métricas de latência e disponibilidade por serviço, e alertas
%     automáticos para a equipa em caso de degradação.
%   - Realizar testes de carga para determinar os limites de capacidade da
%     infraestrutura atual e definir critérios de escalamento horizontal.
%   - Evoluir a estratégia de testes para cobrir os fluxos de streaming do
%     microserviço de IA com fixtures de respostas pré-gravadas.
%   - Integração com sistemas de gestão de lares (ERP) existentes nas instituições,
%     para eliminar duplicação de dados administrativos e facilitar a adoção.

%----------------------------------------------------------------------------------------
\section{Apreciação Final}
\label{sec:apreciacao_final}

% Reflexão pessoal. Pontos a abordar:
%   - O que foi mais desafiante no eixo de backend: provavelmente a implementação
%     do RBAC com isolamento multi-tenant, ou a gestão da ordem de operações no
%     deploy (migrações antes de aceitar tráfego, blue-green sem downtime).
%   - O que foi mais enriquecedor: trabalhar num produto em produção com utilizadores
%     reais impõe uma responsabilidade diferente do ambiente académico — cada decisão
%     de arquitetura tem consequências concretas.
%   - A aprendizagem de que testes e infraestrutura não são atividades separadas do
%     desenvolvimento: os testes de configuração de deploy e a automatização das
%     migrações são código tanto quanto os endpoints da API.
%   - Competências transversais: diagnosticar problemas em produção com informação
%     limitada, comunicar decisões técnicas (e as suas trocas) a uma equipa pequena,
%     e priorizar trabalho num contexto em que tudo parece urgente.
